% !TeX spellcheck = es_ES
%\documentclass[letterpaper,12pt]{article}
\documentclass[letterpaper,11pt]{article}

% Language & Input
  \usepackage[english]{babel} % Language: English
  \usepackage[ansinew]{inputenc} % Input
% Font
  \usepackage[T1]{fontenc} % Font encoding
  \usepackage{lmodern,microtype} % Typeface
% Math
  \usepackage{amsmath,amsthm,amssymb,amsfonts} % AMS math
  \usepackage{dsfont,mathrsfs,ushort} % Math style
% Style
  \usepackage[bf,sf,tiny,center]{titlesec} % Section titles 
  \usepackage{geometry} % Page & margins
  \usepackage{setspace} % Spacing
  \usepackage{enumitem} % Lists
  \usepackage{abstract} % Abstract
  \usepackage{color}
% References
  \usepackage[longnamesfirst]{natbib} % Manages bibiography
% Figures
  \usepackage{graphicx} % Additional options for includefigure
  \usepackage[font={small}]{caption}
% Tables
  \usepackage{booktabs} % Table format
  \usepackage{dcolumn} % Align on the decimal point of numbers in tabular columns.
  \usepackage{multirow} % Multirows in a table
  \usepackage{lscape} % Tables in landscape form
  \usepackage{ctable} % Tables with captions and notes all the same width 
% Hyperlinks
  \usepackage{hyperref}
% Addons
  \usepackage[l2tabu, orthodox]{nag} % Warn the user about the use of such obso­lete things

%%%%%%%%%%%%%%%%%%%%%%%%%%%%%%% Style & Structure
%% Titles (Carlos Perez)
%  \titlelabel{\thetitle.}
%  \titleformat*{\subsection}{\centering\sffamily\itshape}
%  \titleformat{\subsubsection}[runin]{\sffamily}{}{0em}{}[.]
%  \titlespacing{\subsubsection}{0pt}{*2}{\wordsep}
  
% Title page (Luis Cabral)
\makeatletter
    \renewcommand{\@maketitle}{%
    \newpage\parindent0in\null\vskip2em%
    {\LARGE\textbf{\textsf{\@title}}\par\vskip2em}%
    \@author\par\vskip2em%
    \@date%
    \par
    \thispagestyle{empty}\setcounter{page}{1}%\newpage
    }%
\makeatother
\renewenvironment{abstract}% Change the way abstract is presented
{\vskip2em\noindent{\textbf{\textsf{Abstract.}}}\hspace*{1ex}}{}

% Sections
  \renewcommand\thesection{\arabic{section}}
  \renewcommand\thesubsection{\Alph{subsection}}
% Bibliography
  \bibliographystyle{apa} %Alternative: {plain} {aea} {econometrica}
% Spacing & lists and margins
  \geometry{margin=2.54cm}
  \setstretch{1.15} % Set spacing
% Color
  \definecolor{ntr}{rgb}{0.8,0,0} % not too red
% Abstract
  \renewcommand{\abstractnamefont}{\sffamily\small\bfseries}
  \setlength{\abstitleskip}{-1em}
% Hyperlinks
  \hypersetup{colorlinks=true,linkcolor=ntr,citecolor=ntr,urlcolor=black,bookmarksnumbered=true}
% Notes to tables and figures
% Significant stars in tables
  %\newcommand{\sym}[1]{\rlap{#1}}
  \def\sym#1{\ifmmode^{#1}\else\(^{#1}\)\fi}
% Games
% Miscellaneous

%%%%%%%%%%%%%%%%%%%%%%%%%%%%%%% Theorems
  \theoremstyle{plain}
  \newtheorem{prop}{Proposition}
  \newtheorem{assumption} {Assumption}
  \newtheorem{theorem}{Theorem}
  \newtheorem{lem}{Lemma}
  \newtheorem{corollary}{Corollary}
  \theoremstyle{definition}
  \newtheorem{defi}{Definition}
  \newtheorem{ex}{Example}
  \theoremstyle{remark}
  \newtheorem{remark}{Remark}
  \def\propautorefname{Proposition}
  \def\propositionautorefname{Proposition}
  \def\defiautorefname{Definition}
  \def\lemautorefname{Lemma}
  \def\exampleautorefname{Example}
  \def\corollaryautorefname{Corollary}
  \def\exautorefname{Example}

\newcommand{\SR}{\mathbb{R}}

\DeclareMathOperator*{\Max}{m{a}x}
\DeclareMathOperator*{\Min}{m{i}n}
\DeclareMathOperator{\rank}{rank}
\DeclareMathOperator{\plim}{plim}
\DeclareMathOperator{\Var}{Var}
\DeclareMathOperator{\Avar}{AVar}
\DeclareMathOperator{\Cov}{Cov}
\DeclareMathOperator{\E}{\mathbb{E}}
\newcommand{\pto}{\stackrel{p}{\longrightarrow}}
\newcommand{\dto}{\stackrel{d}{\longrightarrow}}
\newcommand{\adistr}{\stackrel{a}{\sim}}
\newcommand{\epsi}{\varepsilon}
\newcommand{\tmle}{\hat \theta_{ML}}
\newcommand{\tgmm}{\hat \theta_{GMM}}
\newcommand\independent{\protect\mathpalette{\protect\independenT}{\perp}}
\def\independenT#1#2{\mathrel{\rlap{$#1#2$}\mkern2mu{#1#2}}}
\DeclareMathOperator{\N}{Normal}

\begin{document}
	
	\title {Econometr\'ia I  \\ Tarea 1}
	\author{Mag\'ister en Econom\'ia, Universidad Alberto Hurtado}
	\date{Fecha de entrega: 10 de septiembre de 2026 (al principio de la ayudant\'ia)}
\maketitle

Total: 90 puntos

\begin{enumerate}
    \item (10 puntos) Suponga un modelo de regresi\'on simple
        \begin{align*}
            y = \beta_0 + \beta_1 x_{1} + u.
        \end{align*}
    Reemplace en la expresi\'on general $\beta=\E(xx')^{-1}\E(xy)$ para demostrar que $\beta_1=\Cov(x_{1},y)/\Var(x_{1})$ y $\beta_0=\E(y)-\beta_1 \E(x_{1})$.
    \item (10 puntos) Estamos interesados en el efecto de una variable $x$ en $y$. Suponga que $x$ es una variable que toma los valores 0, 1 y 2. Defina las siguientes variables aleatorias
    \begin{align*}
    d_1 = 
    \begin{cases}
    1 & \mbox{si } x = 1, \\
    0 & \mbox{en caso contrario, }
    \end{cases}
    d_2 = 
    \begin{cases}
    1 & \mbox{si } x = 2, \\
    0 & \mbox{en caso contrario. }
    \end{cases}
    \end{align*}
    Demuestre que la esperanza condicional $\E(y|d_1,d_2)$ es lineal. \\
    Dado lo visto en clase usted puede deducir que la proyecci\'on lineal es igual a la esperanza condicional. \\ 
    El modelo anterior pertenece a una familia de modelos con la caracter\'istica que la esperanza condicional es lineal ?`Cu\'al es \'esta familia de modelos?
    \item (10 puntos) Utilizando el modelo lineal de regresi\'on $y_i = x_{i}'\beta +u_i$, demuestre que la varianza de $y_i$ se puede descomponer como $\Var(y_i)= \Var(x_i'\beta)+\E(u_i^2)$.\\
    Utilice el resultado anterior para demostrar que $0 \leq R^2 \leq 1$ donde  $R^2=\Var(x_i'\beta)/\Var(y_i)$.
    \item (10 puntos) Suponga dos modelos de regresi\'on 
    \begin{align}
    y &= x'\gamma+u,\\
    y &= x'\beta_1+z'\beta_2+v,
    \end{align}
    donde el modelo (2) incluye $x$ (las variables explicativas del modelo 1) e incluye $z$ (variables adicionales no incluidas en el modelo 1). Adem\'as suponga que $x$ incluye una constante. Demuestre que $\Var(u)\geq \Var(v)$.\\
    Utilice el resultado anterior para argumentar que siempre que agreguemos variables el $R^2$ del modelo no disminuye.
    \item (20 puntos) Dados los modelos de regresi\'on 
    \begin{align}
        y &= \gamma_0+\gamma_1 x_{1}+u,\\
        y &= \beta_0 + \beta_1 x_{1}+\beta_2 x_{2}+v.
    \end{align}
    \begin{enumerate}
        \item Encuentre la relaci\'on entre $\beta_1$ y $\gamma_1$. ?`Cuando se cumple que $\beta_1=\gamma_1$?
        \item Suponga que $y = \log(\text{salarios})$, $x_{1}=$ a\~nos de educaci\'on, y $x_{2}=$ habilidad innata. Si el modelo (4) mide el efecto causal de educaci\'on en salarios, comente sobre las consecuencias de utilizar el modelo (3) en lugar del modelo (4).
    \end{enumerate}
    Pista: Escriba $\gamma_1$ utilizando la f\'ormula para el modelo de regresi\'on simple, y luego reemplace $y$ seg\'un el modelo (4).
    \item (20 puntos) Estamos interesados en el modelo de regresi\'on 
    \begin{align*}
        q^* &= \beta_0 + \beta_1 x_{1}+r^*
    \end{align*}
    donde $\E(r^*)=0$ y $\E(x_1\,r^*)=0$.
    Suponga que no se puede observar $q^*$ pero observamos $q$ tal que $q=q^*+e$ donde $\E(e)=0$ y $\E(x_1\,e)=0$. Se puede interpretar que $q$ es una medici\'on con errores de medida de $q^*$.
    \begin{enumerate}
        \item Dado el modelo lineal de regresi\'on de $q$ (la variable observada con errores de medida) en $x_1$: 
    \begin{align*}
        q &= \delta_0 + \delta_1 x_{1}+error.
    \end{align*}
        Demuestre que $\delta_1=\beta_1$ y $\delta_0=\beta_0$.
        \item Suponga que $\E(r^* e)=0$ y demuestre que la varianza del error es mayor en el modelo con error de medida.\\
    \end{enumerate}
    \item (10 puntos) Dado el proyector lineal de $y$ en $x$, $y=x'\beta+u$ donde $\E(xu)=0$, y el proyector lineal de $\E(y|x)$ en $x$ , $\E(y|x)=x'\beta^*+u^*$ donde $\E(xu^*)=0$. Demuestre que ambos proyectores lineales son id\'enticos, es decir $\beta = \beta^*$.
\end{enumerate}
\end{document} 